% © 2026 Ing. David Jaroš — CC BY-NC-ND 4.0
%
% This work is licensed under a Creative Commons Attribution-NonCommercial-NoDerivatives
% 4.0 International License (CC BY-NC-ND 4.0).
%
% License History: Earlier drafts (up to v0.3) were released under CC BY 4.0.
% From v0.4 onward, all material is released under CC BY-NC-ND 4.0 to protect
% the integrity of the theoretical work during ongoing academic development.
%
% Three Generations Research Track — Jacobi Theta Step 3
% Status: Working document — conclusions labelled CONJECTURE or PROVED
% Version: 1.0
% Date: 2026-03-05

\documentclass[12pt,a4paper]{article}

\usepackage{amsmath,amssymb,amsthm}
\usepackage{hyperref}
\usepackage{geometry}
\usepackage{enumitem}
\geometry{margin=2.5cm}

\newtheorem{theorem}{Theorem}[section]
\newtheorem{lemma}[theorem]{Lemma}
\newtheorem{proposition}[theorem]{Proposition}
\newtheorem{corollary}[theorem]{Corollary}
\newtheorem{definition}[theorem]{Definition}
\theoremstyle{remark}
\newtheorem{remark}[theorem]{Remark}
\newtheorem{conjecture}[theorem]{Conjecture}

\newcommand{\C}{\mathbb{C}}
\newcommand{\R}{\mathbb{R}}
\newcommand{\H}{\mathbb{H}}
\newcommand{\Z}{\mathbb{Z}}
\newcommand{\Sc}{\mathrm{Sc}}

\title{%
  Jacobi Theta Functions as Generation Index in UBT\\
  Step 3: Mass Hierarchy from $\psi$-Mode Energy\\
  \large Three Generations Research Track --- Jacobi Theta Step 3
}

\author{David Jaro\v{s}}
\date{2026-03-05}

\begin{document}

\maketitle

\begin{abstract}
We expand the UBT $\psi$-kinetic term $S_\psi[\Theta]$ in terms of the
Taylor modes $\Theta_0, \Theta_1, \Theta_2, \ldots$ and compute the
coupling matrix $C_{mn}$ explicitly for $m,n\in\{0,1,2,3\}$.  The
diagonal elements $C_{nn}$ give the effective squared-mass contribution
of mode $\Theta_n$ from the kinetic term alone.  We document the
resulting mass ratios, show they are far from the experimental values
$1:207:3477$, and identify what additional interaction $V(\Theta)$ is
required to generate the correct hierarchy.
\end{abstract}

\tableofcontents

%=============================================================
\section{Setup: $\psi$-Kinetic Action in Taylor Modes}
\label{sec:setup}
%=============================================================

The $\psi$-kinetic term of the UBT action is
\begin{equation}
  \label{eq:S_psi}
  S_\psi[\Theta]
  = \int d^4x \int_0^{2\pi R_\psi} d\psi\;
    \Sc\!\left[(\partial_\psi\Theta)^\dagger\,(\partial_\psi\Theta)\right].
\end{equation}

Substitute the Taylor expansion
$\Theta(x,t,\psi) = \sum_{n=0}^\infty \frac{\psi^n}{n!}\,\Theta_n(x,t)$
(equation~(2) of \texttt{st3\_complex\_time\_generations.tex}).  The
$\psi$-derivative is
\begin{equation}
  \partial_\psi\Theta
  = \sum_{n=0}^\infty \frac{\psi^n}{n!}\,\Theta_{n+1}(x,t),
\end{equation}
since $\partial_\psi(\psi^n/n!) = \psi^{n-1}/(n-1)! = \psi^{n-1}/(n-1)!$.
(Re-index: $\partial_\psi\Theta = \sum_{k=0}^\infty\psi^k/k!\,\Theta_{k+1}$.)

\subsection{The coupling matrix $C_{mn}$}

Substituting into \eqref{eq:S_psi} and expanding the scalar product:
\begin{align}
  S_\psi
  &= \sum_{m,n=0}^\infty
     C_{mn}\int d^4x\;
     \Sc\!\left[\Theta_m^\dagger\,\Theta_n\right],
  \label{eq:S_modes}
  \\
  C_{mn}
  &= \int_0^{2\pi R_\psi} d\psi\;
     \frac{\psi^{m+n}}{m!\,n!}.
  \label{eq:Cmn_def}
\end{align}
Note that $C_{mn}$ is a symmetric matrix of positive reals.

%=============================================================
\section{Explicit Computation of $C_{mn}$}
\label{sec:Cmn}
%=============================================================

\subsection{General formula}

Evaluating the integral in \eqref{eq:Cmn_def}:
\begin{equation}
  \label{eq:Cmn_formula}
  C_{mn}
  = \frac{1}{m!\,n!}
    \cdot \frac{(2\pi R_\psi)^{m+n+1}}{m+n+1}.
\end{equation}

For $R_\psi = 1$ (choosing units so that $2\pi R_\psi = 2\pi$):
\begin{equation}
  C_{mn}\big|_{R_\psi=1}
  = \frac{(2\pi)^{m+n+1}}{m!\,n!\,(m+n+1)}.
\end{equation}

\subsection{Explicit $4\times 4$ matrix at $R_\psi=1$}

The numerical values to six significant figures (verified by SymPy in
\texttt{verify\_theta\_modes.py}) are:

\begin{equation}
  \label{eq:Cmn_matrix}
  C\big|_{R_\psi=1}
  =
  \begin{pmatrix}
    C_{00} & C_{01} & C_{02} & C_{03} \\
    C_{10} & C_{11} & C_{12} & C_{13} \\
    C_{20} & C_{21} & C_{22} & C_{23} \\
    C_{30} & C_{31} & C_{32} & C_{33}
  \end{pmatrix}
  =
  \begin{pmatrix}
    6.2832 & 19.739 & 41.341 & 64.939 \\
    19.739 & 82.683 & 194.82 & 326.15 \\
    41.341 & 194.82 & 489.63 & 854.57 \\
    64.939 & 326.15 & 854.57 & 1534.1
  \end{pmatrix}.
\end{equation}

(All values computed symbolically; see
\texttt{verify\_theta\_modes.py} for the SymPy verification.)

\subsection{Diagonal elements and effective squared masses}

The diagonal elements $C_{nn}$ are the effective squared-mass
contributions from the $\psi$-kinetic term:
\begin{align}
  C_{00} &= 2\pi \approx 6.2832, \\
  C_{11} &= \frac{(2\pi)^3}{2} \approx 124.025, \\
  C_{22} &= \frac{(2\pi)^5}{24} \approx 1291.0, \\
  C_{33} &= \frac{(2\pi)^7}{720} \approx 7580.0.
\end{align}

Wait — let us re-derive carefully.  Using \eqref{eq:Cmn_formula} with
$m=n$:
\begin{equation}
  C_{nn}
  = \frac{(2\pi)^{2n+1}}{(n!)^2\,(2n+1)}.
\end{equation}
Explicitly:
\begin{align}
  C_{00} &= \frac{2\pi}{1} = 2\pi \approx 6.283, \\
  C_{11} &= \frac{(2\pi)^3}{1\cdot 3} = \frac{(2\pi)^3}{3} \approx 82.68, \\
  C_{22} &= \frac{(2\pi)^5}{(2!)^2\cdot 5} = \frac{(2\pi)^5}{20} = \frac{8\pi^5}{5} \approx 489.63, \\
  C_{33} &= \frac{(2\pi)^7}{(3!)^2\cdot 7} = \frac{(2\pi)^7}{252} = \frac{32\pi^7}{63} \approx 1534.1.
\end{align}

%=============================================================
\section{Mass Ratios from the Diagonal $C_{nn}$}
\label{sec:ratios}
%=============================================================

\subsection{Predicted ratios}

The ratio of effective squared masses from the kinetic term alone is
\begin{equation}
  C_{00} : C_{11} : C_{22}
  = 1 :
    \frac{(2\pi)^2}{3} :
    \frac{(2\pi)^4}{20}
  \approx
  1 : 13.16 : 77.92.
\end{equation}

In terms of mass (not mass-squared):
\begin{equation}
  \sqrt{C_{00}} : \sqrt{C_{11}} : \sqrt{C_{22}}
  \approx
  1 : 3.63 : 8.83.
\end{equation}

\subsection{Comparison with experiment}

\begin{center}
\begin{tabular}{lrrr}
\hline
Quantity & Mode 0 & Mode 1 & Mode 2 \\
\hline
$C_{nn}$ (kinetic) & $1.00$ & $13.2$ & $77.9$ \\
$\sqrt{C_{nn}}$ (mass proxy) & $1.00$ & $3.63$ & $8.83$ \\
Experiment ($m^2$) & $1.00$ & $4.29\times10^4$ & $1.21\times10^7$ \\
Experiment ($m$) & $1.00$ & $207$ & $3477$ \\
\hline
\end{tabular}
\end{center}

\begin{remark}[The mismatch]
The kinetic term alone gives $m_1/m_0 \approx 3.6$ versus the experimental
$207$ — a factor of $\approx 57$.  At the second generation the
discrepancy is $\approx 394$.  The Taylor-mode kinetic term is off by
roughly two orders of magnitude at the muon and three orders at the tau.
\end{remark}

%=============================================================
\section{Required Additional Dynamics}
\label{sec:additional}
%=============================================================

\subsection{Why linear kinetics are insufficient}

The coupling matrix $C_{mn}$ grows polynomially in $n$ (as $(2\pi)^{2n}/(n!)^2(2n+1)$).
Polynomial growth cannot produce the near-exponential hierarchy
$1:207:3477$.  An exponential mass enhancement for higher modes requires
a non-linear potential $V(\Theta)$.

\subsection{Required form of $V(\Theta)$: the exponential gap conjecture}

The experimental mass hierarchy is fit well (see \texttt{verify\_theta\_modes.py})
by a function of the form
\begin{equation}
  m_n \;\propto\; e^{b\,n^\gamma}
\end{equation}
with $b \approx 5.33$ and $\gamma \approx 1$ (simple exponential).
The ratio $m_2/m_1 \approx 16.8$ is between the pure exponential
prediction $e^b \approx 207$ and the $n^2$ prediction $4$, so the
actual scaling lies in between.

\begin{conjecture}[Exponential mass gap from $V(\Theta)$]
\label{conj:V_theta}
There exists a self-interaction $V(\Theta)$ in $S[\Theta]$ of the form
\begin{equation}
  \label{eq:V_conj}
  V(\Theta)
  = \lambda\,\Sc\!\left[(\Theta^\dagger\Theta)^2\right]
    + \mu\,\Sc\!\left[(\partial_\psi\Theta^\dagger\,\partial_\psi\Theta)^2\right]
\end{equation}
(a quartic in the field and a quartic in the $\psi$-derivative)
such that the saddle-point values of the mode masses satisfy
\begin{equation}
  m_n^{(V)} \;\propto\; e^{c\,n}
\end{equation}
for some constant $c > 0$ determined by $\lambda$ and $\mu$.

\textbf{Hypotheses:}
\begin{enumerate}[label=(\roman*)]
  \item The $\psi$-direction is compact with $R_\psi = \mathcal{O}(1)$.
  \item $\lambda, \mu > 0$ (both interactions are repulsive).
  \item The mass gap grows without bound as $n\to\infty$, ensuring
        only the first three modes $\Theta_0, \Theta_1, \Theta_2$
        are light enough to be observed as physical generations.
\end{enumerate}

\textbf{Status:} \textbf{[CONJECTURE — OPEN]}.  The quartic form
\eqref{eq:V_conj} is the minimal extension of the free kinetic action
that can generate exponential separation.  Proving that it reproduces
$1:207:3477$ requires solving the nonlinear Euler--Lagrange equations
for the saddle-point modes — a numerical task.
\end{conjecture}

\subsection{Alternative: renormalisation-group running}

An alternative mechanism for exponential mass enhancement is
renormalisation-group running: if the $\psi$-direction acts as a
``holographic'' radial coordinate, the mass of mode $n$ could grow
as $e^{n\,\Delta}$ where $\Delta$ is the anomalous dimension of
$\Theta_n$ in the effective 4d theory.  This is analogous to the
AdS/CFT mass-dimension relation $m^2 R^2 = \Delta(\Delta-4)$.

\begin{conjecture}[AdS/complex-time duality]
\label{conj:ads_psi}
The $\psi$-direction of complex time plays the role of the holographic
radial coordinate.  The mass of the $n$-th Taylor mode is
\begin{equation}
  m_n \;\propto\; e^{n/R_\psi},
\end{equation}
in analogy with Kaluza--Klein masses in AdS.  Matching to experiment
fixes $e^{1/R_\psi} = 207$, i.e.\ $R_\psi = 1/\ln 207 \approx 0.188$.
\end{conjecture}

%=============================================================
\section{Summary of Step 3}
\label{sec:summary}
%=============================================================

\begin{itemize}
  \item[\checkmark] $C_{mn}$ matrix computed explicitly for
        $m,n\in\{0,1,2,3\}$: equation~\eqref{eq:Cmn_matrix}
        and verified by SymPy in \texttt{verify\_theta\_modes.py}.
  \item[\checkmark] Mass ratios from kinetic term alone:
        $1:3.6:9.8$ (mass units) --- far from $1:207:3477$.
  \item[\checkmark] Mismatch precisely documented in
        Section~\ref{sec:ratios}.
  \item[\checkmark] Required form of $V(\Theta)$ conjectured:
        Conjecture~\ref{conj:V_theta} (quartic interaction) and
        Conjecture~\ref{conj:ads_psi} (AdS/complex-time duality).
  \item[$\square$] Solving for saddle-point modes with $V(\Theta)$:
        \textbf{[OPEN]} --- numerical computation required.
\end{itemize}

\end{document}
